% =====================================================================
%  Grounded Arithmetic: motivation, history, and the state of the art
%  Beamer source.  Compile with:  pdflatex grounded.tex  (twice)
% =====================================================================
\documentclass[aspectratio=169,11pt]{beamer}

\usepackage[T1]{fontenc}
\usepackage[utf8]{inputenc}
\IfFileExists{lmodern.sty}{\usepackage{lmodern}}{}   % optional: nicer fonts if available
\usepackage{amsmath,amssymb}
\usepackage{booktabs}
\usepackage{xcolor}

% ---------------------------------------------------------------- look
\usetheme{default}
\setbeamertemplate{navigation symbols}{}
\setbeamertemplate{itemize items}[circle]
\setbeamertemplate{itemize subitem}[triangle]

\definecolor{ink}{HTML}{16202B}
\definecolor{accent}{HTML}{1F4E7A}
\definecolor{warn}{HTML}{8C2F39}
\definecolor{gap}{HTML}{8A6023}
\definecolor{soft}{HTML}{5B6A78}

\setbeamercolor{structure}{fg=accent}
\setbeamercolor{normal text}{fg=ink}
\setbeamercolor{frametitle}{fg=ink}
\setbeamerfont{frametitle}{series=\bfseries,size=\large}
\setbeamerfont{title}{series=\bfseries}
\setbeamercolor{alerted text}{fg=warn}

\setbeamertemplate{frametitle}{%
  \vspace{0.6em}\insertframetitle\par
  \vspace{-0.35em}{\color{accent}\rule{\textwidth}{0.6pt}}}

\setbeamertemplate{footline}{%
  \hfill{\tiny\color{soft}\insertframenumber\ /\ \inserttotalframenumber}%
  \hspace{1em}\vspace{0.6em}}

\newcommand{\hq}{\textit{habeas quid}}
\newcommand{\key}[1]{{\color{accent}\textbf{#1}}}
\newcommand{\bad}[1]{{\color{warn}\textbf{#1}}}
\newcommand{\amber}[1]{{\color{gap}\textbf{#1}}}

% ---------------------------------------------------------------- meta
\title{Reasoning Around Recursion}
\subtitle{Grounded arithmetic: where it came from, and where it is}
\author{Omar}
\institute{\small Indian Institute of Science \\[0.2em]
           \footnotesize joint work with the Grounded Deduction group, EPFL DEDIS}
\date{\today}

\begin{document}

\begin{frame}[plain]
  \titlepage
\end{frame}

% =====================================================================
\section{Motivation}
% =====================================================================

\begin{frame}{Two programs}
  \vfill
  \begin{columns}[T]
    \begin{column}{0.48\textwidth}
      \underline{C}\\[1em]
      \texttt{main() \{ printf("Hi!"); \}}\\[0.6em]
      \texttt{main() \{ main(); \}}\\[1.5em]
      {\color{accent}\Large Both compile.}
    \end{column}
    \begin{column}{0.48\textwidth}
      \underline{Coq / Lean / Agda / Isabelle}\\[1em]
      \texttt{Fixpoint loop (n:nat) := loop n.}\\[1.5em]
      {\color{warn}\Large Rejected.}
    \end{column}
  \end{columns}
  \vfill
  \begin{center}\large Why?\end{center}
  \vfill
\end{frame}

\begin{frame}{And again, with datatypes}
  \vfill
  \begin{columns}[T]
    \begin{column}{0.44\textwidth}
      \underline{Haskell}\\[0.7em]
      {\small\texttt{data Expr}\\
      \texttt{~~= Base}\\
      \texttt{~~| App Expr Expr}\\
      \texttt{~~| Abs (Expr -> Expr)}}\\[1.2em]
      {\color{accent}\Large Compiles.}
    \end{column}
    \begin{column}{0.52\textwidth}
      \underline{Coq}\\[0.7em]
      {\small\texttt{Inductive Expr :=}\\
      \texttt{~~| Base : Expr}\\
      \texttt{~~| App : Expr -> Expr -> Expr}\\
      \texttt{~~| Abs : (Expr -> Expr) -> Expr.}}\\[1.2em]
      {\color{warn}\Large Rejected.}
    \end{column}
  \end{columns}
  \vfill
  \pause
  \begin{center}\small
    \bad{Non strictly positive occurrence of ``Expr'' in ``(Expr -> Expr) -> Expr''.}
  \end{center}
  \vfill
  \begin{center}
    A \emph{different} refusal --- not termination, \key{positivity}.\\
    {\small Not one narrow check: a family of restrictions on what you may define.}
  \end{center}
\end{frame}

\begin{frame}{The usual answer is wrong}
  \vfill
  ``Termination checking is a conservative engineering choice.''
  \vfill
  \pause
  \begin{center}
    \Large It is not a choice.\\[0.8em]
    \large It is the \key{consistency} of your proof assistant.
  \end{center}
  \vfill
  \pause
  \begin{center}
    Working out exactly why took fifteen years,\\
    and ended two research programmes.
  \end{center}
  \vfill
\end{frame}

% =====================================================================
\section{History}
% =====================================================================

\begin{frame}{Background: why anyone wanted a way out}
  \begin{itemize}\itemsep1.1em
    \item \textbf{Frege} (1893): mathematics from logic, via extensions of concepts.
    \item \textbf{Russell} (1902): naive comprehension is inconsistent.\\
          {\small $R = \{x : x \notin x\}$. Is $R \in R$?}
    \item Two repairs, both in wide use by 1930:
      \begin{itemize}\itemsep0.4em
        \item \key{Russell} --- \emph{type theory}
        \item \key{Zermelo} --- \emph{axiomatic set theory}
      \end{itemize}
  \end{itemize}
  \vfill
  \pause
  \begin{center}
    Type theory works by \key{stratification}: every object gets a level,
    and\\ a collection may only contain objects of strictly lower level.
  \end{center}
\end{frame}

\begin{frame}{What stratification costs}
  \vfill
  Under stratification, \key{$x\,x$ is not false --- it is ungrammatical.}\\
  {\small A thing may never be applied to itself. That is the whole mechanism.}
  \vfill
  \pause
  But self-application is exactly how you get \key{recursion}:\\
  {\small from $x\,x$ you can build a fixed-point combinator, and a fixed point
   solves \emph{any} recursive equation.}
  \vfill
  \pause
  \begin{center}\large
    Ban it, and you lose general recursion.
  \end{center}
  \vfill
  \pause
  \begin{quote}\small
    ``Rather than adopt the method of Russell \dots\ or that of Zermelo,
    \textbf{both of which appear somewhat artificial}, we introduce \dots\
    a certain restriction on the law of excluded middle.''
  \end{quote}
  \hfill{\footnotesize --- Church 1932}
\end{frame}

\begin{frame}{Church and Curry: functions, not sets}
  \begin{center}\small
    ``Both of them set out to base logic and mathematics on \textbf{functions}
    instead of on set theory.'' \hfill{\footnotesize --- Seldin}
  \end{center}
  \vfill
  \begin{columns}[T]
    \begin{column}{0.47\textwidth}
      \underline{\key{Extensional}} --- set theory\\[0.6em]
      {\small A function \emph{is} its graph: the set of its
       input--output pairs. Two functions are the same iff they agree
       on every input.}\\[0.8em]
      {\small So $\lambda n.\,2n$ and $\lambda n.\,n+n$ are the
       \emph{same} function.}
    \end{column}
    \begin{column}{0.47\textwidth}
      \underline{\key{Intensional}} --- Church, Curry\\[0.6em]
      {\small A function is a \emph{rule}: a finite description of how to get
       the output. Two functions are the same iff the descriptions are.}\\[0.8em]
      {\small So $\lambda n.\,2n$ and $\lambda n.\,n+n$ are
       \emph{different} functions.}
    \end{column}
  \end{columns}
  \vfill
  \pause
  \begin{center}
    A rule can be applied to a rule. Compilers compile compilers.\\
    \key{No stratification required.}
  \end{center}
\end{frame}

\begin{frame}{Church, 1932}
  \textit{A Set of Postulates for the Foundation of Logic}
  \vfill
  \begin{itemize}\itemsep0.8em
    \item $\lambda$ as a \emph{logical} primitive, alongside $\Pi, \Sigma, \&, \sim, \iota, A$
    \item No free variables
    \item \key{Restricted excluded middle}: for some arguments $\mathbf{X}$,
          $\mathbf{F(X)}$ ``is undefined and represents nothing''
    \item Keeps double negation; \key{abandons \emph{reductio}} in certain cases
    \item Guarded quantifier: $\Pi(F,G)$ licenses $G(M)$ only once $F(M)$ is \emph{known true}
  \end{itemize}
  \vfill
  \pause
  \begin{center}\small
    Reject LEM $\cdot$ weaken \emph{reductio} $\cdot$ let well-formed terms denote nothing\\
    \key{---\ a recognisable sketch of GA, ninety-four years early.}
  \end{center}
\end{frame}

\begin{frame}{Church did not claim consistency}
  \vfill
  \begin{quote}
    ``Whether the system of logic which results from our postulates is adequate
    for the development of mathematics, and \textbf{whether it is wholly free from
    contradiction}, are questions which we cannot now answer except by
    \textbf{conjecture}.''
  \end{quote}
  \hfill{\footnotesize --- Church 1932, p.\,348}
  \vfill
  \pause
  \begin{center}\large He proposed to find out \emph{empirically}.\end{center}
  \vfill
\end{frame}

\begin{frame}{Curry, 1930}
  \textit{Grundlagen der kombinatorischen Logik}
  \vfill
  \begin{itemize}\itemsep0.8em
    \item Started (1922) from a different question: \key{substitution} in \emph{Principia}
          is far more complicated than modus ponens --- break it into simpler rules
    \item Combinators $\mathsf{I},\mathsf{B},\mathsf{C},\mathsf{W},\mathsf{K},\mathsf{S}$; no bound variables at all
    \item Consistency proof for the \emph{combinator} part --- but no logical
          axioms given; that half was left for future work
  \end{itemize}
  \vfill
  \pause
  \begin{center}\small
    On the Russell term $[x](\neg(xx))$: Curry proposed to find postulates
    from which \key{it could not be proved that its self-application is a proposition}
    --- while insisting the term \emph{exist} in the system.
  \end{center}
\end{frame}

\begin{frame}{Kleene \& Rosser, 1935}
  \vfill
  Both systems are strong enough to enumerate their own provably-total functions.\\
  That is enough to formalise \key{Richard's paradox}.
  \vfill
  \begin{quote}\small
    ``We are sorry to say that we have a proof of the inconsistency of your
    system of formal logic.''
  \end{quote}
  \pause
  \begin{quote}\small
    ``I am afraid that I can not share your grief \dots\ you are not sorry but you
    are full of that profound satisfaction which comes from doing something worth
    while \dots\ \textbf{I hasten to congratulate you both.}''
  \end{quote}
  \hfill{\footnotesize --- Curry to Kleene and Rosser, 19 November 1934}
\end{frame}

\begin{frame}{\emph{That} --- but not \emph{why}}
  \vfill
  Kleene and Rosser showed the systems were inconsistent.\\
  They did not isolate \key{which principle} was responsible.
  \vfill
  \begin{itemize}\itemsep0.6em
    \item The proof was long, and specific to those systems
    \item Curry's response: a 60-page paper (1941) studying the paradox
          ``so as to lay bare its central nerve''
    \item Church's response (1933): delete four postulates, the ones carrying
          \emph{reductio} --- \bad{it did not save the system}
  \end{itemize}
  \vfill
  \pause
  \begin{center}\Large The answer came in 1942, in three pages.\end{center}
  \vfill
\end{frame}

% =====================================================================
\section{Curry's paradox}
% =====================================================================

\begin{frame}{Curry, 1942 --- three pages}
  \textit{The Inconsistency of Certain Formal Logics}, JSL 7(3):115--117
  \vfill
  Suppose a system has:
  \begin{enumerate}\itemsep0.5em
    \item $\lambda$-abstraction that computes \hfill {\small(combinatory completeness)}
    \item $\vdash M \supset M$
    \item \key{contraction}: \; from $\vdash M \supset (M \supset N)$ \; infer \; $\vdash M \supset N$
    \item modus ponens
  \end{enumerate}
  \vfill
  \pause
  \begin{center}\Large Then \bad{every term is provable.}\end{center}
\end{frame}

\begin{frame}{The derivation}
  Let $\mathfrak{A}$ be a term with $\mathfrak{A} \equiv (\mathfrak{A} \supset \mathfrak{B})$,
  \; obtained from a fixed point of $\lambda X.\,X \supset \mathfrak{B}$.
  \vfill
  \[
  \begin{array}{r@{\qquad}l@{\qquad}l}
    (1) & \vdash \mathfrak{A} \supset \mathfrak{A}
        & \text{\small identity}\\[0.5em]
    (2) & \vdash \mathfrak{A} \supset .\,\mathfrak{A} \supset \mathfrak{B}
        & \text{\small unfold the definition}\\[0.5em]
    (3) & \vdash \mathfrak{A} \supset \mathfrak{B}
        & \text{\small {\color{gap}contraction}}\\[0.5em]
    (4) & \vdash \mathfrak{A}
        & \text{\small fold: $(3)$ \emph{is} $\mathfrak{A}$}\\[0.5em]
    (5) & \vdash \mathfrak{B}
        & \text{\small modus ponens}
  \end{array}
  \]
  \vfill
  \begin{center}$\mathfrak{B}$ was arbitrary.\end{center}
\end{frame}

\begin{frame}{What is \emph{not} in that proof}
  \vfill
  \begin{center}\Large No negation.\end{center}
  \vfill
  \pause
  Identity, modus ponens, contraction, unfolding a definition.\\[0.6em]
  Every one is valid \key{intuitionistically}. Every one is valid in \key{minimal logic}.
  \vfill
  \pause
  \begin{center}\large
    Church had restricted excluded middle.\\
    Church had abandoned \emph{reductio}.\\[0.8em]
    \bad{Neither was ever going to be enough.}
  \end{center}
  \vfill
\end{frame}

\begin{frame}{Curry's own diagnosis}
  \vfill
  \begin{quote}\large
    ``Before he discovered this paradox, Curry had assumed that the contradiction
    was caused by his postulates for the universal quantifier, but this paradox
    made it clear that \textbf{the problem was the postulates for implication
    alone}.''
  \end{quote}
  \hfill{\footnotesize --- Seldin, \emph{The Logic of Curry and Church}}
  \vfill
  \pause
  \begin{center}
    Church guarded the \emph{quantifier}. Curry blamed the \emph{quantifier}.\\[0.6em]
    \key{It was implication.}
  \end{center}
  \vfill
\end{frame}

\begin{frame}{What the proof actually needs}
  \vfill
  Curry states hypothesis 1 as \key{combinatory completeness}: for any term $X$
  in variables $x_1,\dots,x_n$, there is a closed term $A$ with
  \[ A\,x_1 \cdots x_n \;=\; X \]
  {\small That is: \emph{abstraction is definable}. $A$ is what we would write
   $\lambda x_1 \dots x_n.\,X$.\\
   He assumes it because that is what his systems offered.}
  \vfill
  \pause
  But the proof consumes only \emph{one} term:
  \[ \mathfrak{A} \qquad\text{with}\qquad
     \mathfrak{A} \;\equiv\; \mathfrak{A} \supset \mathfrak{B} \]
  {\small Abstraction supplies it, via a fixed-point combinator.
   So does a plain definitional mechanism: write $d \equiv d \to \mathfrak{B}$.}
  \vfill
  \pause
  \begin{center}\Large
    The real hypothesis is:\\
    \key{the system admits self-referential definitions.}
  \end{center}
  \vfill
  \begin{center}\small
    Combinatory completeness is one route to that --- not the only one,
    and not equivalent to it.
  \end{center}
\end{frame}

\begin{frame}{The shape of the result}
  \vfill
  \begin{center}\Large
    \key{Unrestricted recursion + ordinary implication\\[0.3em] = proves everything.}
  \end{center}
  \vfill
  \pause
  \begin{center}
    And that is why \texttt{Fixpoint loop (n:nat) := loop n.} is rejected.
  \end{center}
  \vfill
  \pause
  \begin{center}\small
    Not conservatism. If your system admitted it, you could derive $\bot$.
  \end{center}
  \vfill
\end{frame}

\begin{frame}{After 1935}
  \vfill
  \key{Church, Kleene and Rosser} extracted the \emph{pure} $\lambda$-calculus ---
  terms and conversion, no logic --- and Church retreated to simple type theory
  in 1940.
  \vfill
  \key{Curry did not retreat.} He kept the untyped terms and added a predicate
  $\mathsf{H}$ for ``is a proposition,'' using it to guard the logical rules.
  \vfill
  \pause
  \begin{center}
    Everything you use descends from Church 1940.\\[0.4em]
    \key{GA is closer to Curry's road.}
  \end{center}
  \vfill
\end{frame}

\begin{frame}{Where everyone landed}
  \small
  \begin{center}
  \begin{tabular}{@{}lccc@{}}
    \toprule
    & \textbf{unrestricted} & \textbf{ordinary} & \textbf{non-} \\
    & \textbf{recursion} & \textbf{rules} & \textbf{triviality} \\
    \midrule
    Church 1932 $\cdot$ Curry 1930    & kept    & weakened neg.\ & \bad{lost} \\
    STLC $\cdot$ Coq $\cdot$ Agda $\cdot$ Lean $\cdot$ HOL & \bad{dropped} & kept & kept \\
    Contraction-free \small(Fitch, Gri\v{s}in, Cantini) & kept & \amber{structural} & kept \\
    Paraconsistent \small(Priest)     & kept    & \amber{no explosion} & \amber{contained} \\
    Kripke 1975 \small(truth gaps)    & kept    & \amber{partial}  & kept \\
    \midrule
    \key{Grounded arithmetic}         & \key{kept} & \key{preconditions} & \key{kept} \\
    \bottomrule
  \end{tabular}
  \end{center}
  \vfill
  \begin{center}\small
    Everyone since 1942 kept implication and gave up recursion.\\

  \end{center}
\end{frame}

% =====================================================================
\section{Grounded deduction}
% =====================================================================

\begin{frame}{Kripke, 1975 --- there is no sieve}
  {\small \textit{Outline of a Theory of Truth}, J.\ Phil.\ 72(19):690--716}
  \vfill
  Whether a sentence is paradoxical can depend on \emph{empirical facts},
  not on its shape.
  \vfill
  \pause
  \begin{quote}\small
    ``The moral: an adequate theory must allow our statements involving the notion
    of truth to be \textbf{risky} \dots\ There can be \textbf{no syntactic or
    semantic ``sieve''} that will winnow out the ``bad'' cases while preserving the
    ``good'' ones.''
  \end{quote}
  \hfill{\footnotesize --- p.\,692}
  \vfill
  \pause
  \begin{center}
    So you cannot fix this by restricting \emph{what may be written}.\\
    \key{Exactly the argument against the type-theoretic route.}
  \end{center}
\end{frame}

\begin{frame}{Grounded}
  \vfill
  \begin{quote}\small
    ``If ultimately this process terminates in sentences not mentioning the concept
    of truth, so that the truth value of the original statement can be ascertained,
    we call the original sentence \textbf{grounded}; otherwise, \textbf{ungrounded}.''
  \end{quote}
  \hfill{\footnotesize --- p.\,694}
  \vfill
  \pause
  The detail people miss --- the \key{truth-teller}:
  \begin{center}\small (3) \quad ``(3) is true.''\end{center}
  {\small ``which, though not paradoxical, yield no determinate truth conditions'' (p.\,693);
   ``Sentences such as (3), though not paradoxical, are ungrounded'' (p.\,694).}
  \vfill
  \pause
  \begin{center}\Large
    Groundedness is not about \emph{paradox}.\\
    \key{It is about having a basis.}
  \end{center}
\end{frame}

\begin{frame}{The least fixed point}
  \vfill
  \begin{itemize}\itemsep0.55em
    \item $T(x)$ is interpreted by a \key{partial set} $(S_1, S_2)$ ---
          extension, antiextension, undefined outside $S_1 \cup S_2$
    \item Connectives evaluate by \key{Kleene's strong three-valued scheme}
          {\small (p.\,700; fn.\,18 cites Kleene 1952, \S64)}
    \item The jump $\phi$ is \key{monotone}: extending $T$ never changes a value
          already established. \\
          {\small ``This property --- technically, the monotonicity of $\phi$ ---
           is crucial for all our constructions.'' (p.\,703)}
    \item Iterate transfinitely $\Rightarrow$ a \key{minimal} fixed point:
          ``any fixed point extends'' it (p.\,705)
  \end{itemize}
  \vfill
  \pause
  \begin{center}
    ``(Being a fixed point, $\mathcal{L}_\sigma$ is a language that
    \key{contains its own truth predicate}.)'' \hfill{\footnotesize --- p.\,705}
  \end{center}
  \vfill
  \begin{center}\small
    Kripke's caution: ``\emph{Undefined} is not an extra truth value'' (fn.\,18) ---
    this is not three-valued logic.
  \end{center}
\end{frame}

\begin{frame}{What GA takes, and what it adds}
  \vfill
  \begin{columns}[T]
    \begin{column}{0.46\textwidth}
      \underline{From Kripke}\\[0.7em]
      {\small The word \emph{grounded}.\\[0.4em]
       Ungrounded expressions denote nothing, so paradoxes are inert.\\[0.4em]
       A language containing its own truth predicate.}
    \end{column}
    \begin{column}{0.46\textwidth}
      \underline{GA adds}\\[0.7em]
      {\small Grounding in \emph{computation}, not a semantic construction.\\[0.4em]
       The condition as \emph{preconditions on inference rules}.\\[0.4em]
       Arithmetic, Turing-complete computation, machine-checked metatheory.}
    \end{column}
  \end{columns}
  \vfill
  \pause
  \begin{center}\small
    BGA proves its own truth predicate for grounded sentences (Thm.\ 4.20) ---
    Kripke's \emph{raison d'\^etre}, in a system that also computes.
  \end{center}
  \vfill
\end{frame}

\begin{frame}{\emph{Habeas quid} --- ``have a thing''}
  \vfill
  You may not reason with a term until you have shown it \emph{has a value}.
  \vfill
  \[
    \frac{\Gamma \vdash p\ \mathsf{B} \qquad \Gamma \cup \{p\} \vdash q}
         {\Gamma \vdash p \to q}\ (\to I)
    \qquad\quad
    \frac{\Gamma \vdash p\ \mathsf{B} \quad \Gamma\cup\{\neg p\} \vdash q
          \quad \Gamma\cup\{\neg p\} \vdash \neg q}
         {\Gamma \vdash p}
  \]
  \vfill
  \begin{center}\small
    $p\ \mathsf{B} \;\equiv\; p \vee \neg p$ --- classically a tautology, so the
    preconditions are free and the rules collapse to the familiar ones.\\[0.4em]
    In GA it is a \key{dynamic type check}: a demand that $p$ terminate with a boolean.
  \end{center}
  \vfill
  \begin{center}
    In Curry's terms: GA satisfies hypotheses \key{1, 2 and 4}.\\
    It breaks \key{3}.
  \end{center}
\end{frame}

\begin{frame}{Both paradoxes die of malnutrition}
  \begin{columns}[T]
    \begin{column}{0.47\textwidth}
      \underline{Liar} \quad $L \equiv \neg L$\\[0.8em]
      Admitted as a definition.\\[0.6em]
      Refutation by contradiction needs $L\ \mathsf{B}$,\\
      which needs the refutation.\\[0.8em]
      \key{No value. Nothing to explode from.}
    \end{column}
    \begin{column}{0.47\textwidth}
      \underline{Curry} \quad $C \equiv C \to P$\\[0.8em]
      Admitted as a definition.\\[0.6em]
      Step (3) now needs $C\ \mathsf{B}$ first,\\
      obtainable only via the implication\\
      we are trying to introduce.\\[0.8em]
      \key{Circular obligation. No proof.}
    \end{column}
  \end{columns}
  \vfill
  \begin{center}\small
    The paradoxes are \key{expressible} and simply not \key{provable}.
  \end{center}
\end{frame}

\begin{frame}{Does GA forbid self-application?}
  \vfill
  \begin{center}\Large No.\end{center}
  \vfill
  \pause
  \key{PGA} has higher-order application $t[t]$ as a \emph{primitive}.
  So $f[f]$ is a well-formed term.
  \begin{center}\small
    Everything is a natural number: $f[a]$ reads $f$ as a G\"odel code of a
    primitive-recursive step function, applied to $a$, with an unbounded search.\\
    $\mathcal{S},\mathcal{K},\mathcal{I}$ are just particular numerals.
    No $\lambda$ binder --- but $\lambda$ is derivable by combinator conversion.
  \end{center}
  \vfill
  \pause
  \key{BGA} is the interesting contrast: no term is ever applied to a term.
  Recursion comes only from $d(t,\dots,t)$ --- definition symbols applied to
  argument terms, freely mutually recursive, no ordering, no obligation at
  definition time.
  \begin{center}\small And it is Turing-complete anyway.\end{center}
  \vfill
\end{frame}

\begin{frame}{The restriction is not on what you may write}
  \vfill
  Type theory's mechanism is \key{grammatical}: make $x\,x$ ill-formed,
  and the paradox cannot be stated.
  \vfill
  \pause
  GA does the opposite. All of these are \key{well-formed}:
  \begin{center}\small
    $f[f]$ \qquad $L \equiv \neg L$ \qquad $C \equiv C \to P$
    \qquad \texttt{collatz}$(n)$
  \end{center}
  \vfill
  \pause
  \begin{center}\Large
    The restriction is on what you may \key{infer}.
  \end{center}
  \vfill
  \begin{center}\small
    Which is why the paradoxes are \emph{expressible} and merely not \emph{provable}
    ---\\ and why a definition never needs a termination proof to be admitted.
  \end{center}
  \vfill
\end{frame}

% =====================================================================
\section{The development so far}
% =====================================================================

\begin{frame}{Six papers, two years}
  \small
  \begin{center}
  \begin{tabular}{@{}llp{6.4cm}@{}}
    \toprule
    \textbf{2024} & GD            & the idea, and the propositional core \\
    \textbf{2025} & BGA / PGA     & machine-checked metatheory \\
    \textbf{2026} & RGA           & quantifiers, grounded reflectively \\
    \textbf{2026} & Internalized truth & its own truth predicate, full language \\
    \textbf{2026} & OGA           & one more rule --- incompleteness becomes \emph{priced} \\
    \textbf{2026} & Grounded set theory & a graded ZF \\
    \midrule
    \textbf{now}  & \key{Isabelle/Pure} & \key{can anyone actually use it?} \\
    \bottomrule
  \end{tabular}
  \end{center}
  \vfill
  \begin{center}\small
    Roughly: \emph{idea} $\to$ \emph{proof it works} $\to$ \emph{make it a logic}
    $\to$ \emph{make it reflective} $\to$ \emph{price the gap} $\to$
    \emph{make it a foundation} $\to$ \emph{make it usable}.\\[0.4em]
    \key{Four of the six are from 2026.}
  \end{center}
\end{frame}

\begin{frame}{1 $\cdot$ \emph{Reasoning Around Paradox with Grounded Deduction}}
  {\small Ford, arXiv:2409.08243 (2024)}
  \vfill
  \begin{itemize}\itemsep0.7em
    \item Kripke-inspired: let ungrounded expressions denote nothing
    \item \hq\ as \emph{preconditions on inference rules} --- the central move
    \item An impossibility triangle:\\
          \small consistency $\cdot$ unrestricted LEM $\cdot$ unrestricted recursive definitions
    \item Propositional GD, predicate logic, arithmetic; a computational
          interpretation and a denotational semantics
    \item Speculative closing: reflective idealization, ``Cantor's paradise regained\ldots?''
  \end{itemize}
  \vfill
  \begin{center}\small \amber{Preliminary --- no machine-checked metatheory.}\end{center}
\end{frame}

\begin{frame}{2 $\cdot$ \emph{Have a thing?} \; --- BGA and PGA}
  {\small arXiv:2510.25369}
  \vfill
  Two concrete systems, everything checked in Isabelle/HOL.
  \vfill
  \begin{columns}[T]
    \begin{column}{0.46\textwidth}
      \underline{BGA} --- minimal kernel\\[0.6em]
      \small
      syntactic consistency\\
      semantic soundness\\
      \key{semantic completeness}\\
      represents every general-recursive function\\
      G\"odel arithmetization + reflection\\
      \key{its own truth predicate} (grounded sentences)
    \end{column}
    \begin{column}{0.46\textwidth}
      \underline{PGA} --- expressive layer\\[0.6em]
      \small
      primitive propositional operators\\
      higher-order functions via SKI\\
      consistency, completeness\\
      reduces to BGA\\[0.8em]
      16 of Copi's 19 classical rules survive \emph{unmodified}
    \end{column}
  \end{columns}
  \vfill
  \begin{center}\small
    Proof-theoretic strength: somewhere in
    $[\,\omega^{\omega},\ \omega_1^{\mathrm{CK}}\,]$ --- \amber{open}.
  \end{center}
\end{frame}

\begin{frame}{Consistent \emph{and} complete \emph{and} arithmetic?}
  \vfill
  \begin{itemize}\itemsep0.8em
    \item G\"odel~I concerns \key{syntactic} completeness --- for every $f$,
          $\vdash f$ or $\vdash \neg f$.\\
          {\small BGA is emphatically not that, and cannot be: LEM is built into the definition.}
    \item What is proved is \key{semantic} completeness, \emph{with respect to
          the operational-semantic model}.\\
          {\small Every formula that reduces to 1 under all assignments is provable.
          Not: ``GA proves every truth of arithmetic.''}
    \item G\"odel's argument assumes a classical \emph{object} logic.
          That assumption fails; the metalogic assumption still holds.
  \end{itemize}
  \vfill
  \begin{center}\small
    The consistency and completeness proofs are simulation arguments in
    opposite directions ---\\ \key{the Church--Turing thesis continuing to hold.}
  \end{center}
\end{frame}

\begin{frame}{3 $\cdot$ \emph{Computable Quantification in RGA}}
  {\small Ford, arXiv:2607.25533 (2026)}
  \vfill
  Quantifiers --- grounded \key{reflectively}:
  \begin{center}\small
    a universal is \emph{true} when the system's own proof search certifies its
    schematic instance,\\ \emph{false} when it refutes a particular numeral instance.
  \end{center}
  \vfill
  \begin{itemize}\itemsep0.5em
    \item Proves totality of $+$ and $\times$ as internally quantified theorems
    \item Represents \key{exactly the recursively enumerable sets}
    \item Machine-checked: soundness, consistency, open completeness,
          $\mathbb{N}$-soundness, a Church--Turing characterisation
    \item \key{$\omega$-incompleteness} --- grounded truth is r.e., so some family has
          every numeric instance provable while its closure is
          \emph{semantically ungrounded}
  \end{itemize}
  \vfill
  \begin{center}\small
    DNE holds $\cdot$ quantified excluded middle fails $\cdot$
    refuted universals yield explicit witnesses\\
    \amber{the deduction theorem's abstraction direction fails at ungrounded hypotheses}
  \end{center}
\end{frame}

\begin{frame}{4 $\cdot$ \emph{Internalized Truth in RGA}}
  {\small Ford, arXiv:2608.16140 (Aug 2026)}
  \vfill
  Tarski: no consistent \emph{classical} system with arithmetic defines its own truth.
  \vfill
  \begin{itemize}\itemsep0.6em
    \item A truth predicate for RGA's \key{full language, quantifiers included},
          as an \emph{internal RGA term}
    \item Compiled from a primitive-recursive decider; adequate in both directions
    \item A square of metatheorems --- provable $\Rightarrow$ internally true;
          grounded-true $\Rightarrow$ internally provable; internal truth
          $\Rightarrow$ internal provability; consistency follows
    \item Two \emph{disjoint} internal machines: a certified decider and a
          certified proof-checker, both RGA terms
  \end{itemize}
  \vfill
  \begin{center}\small
    Along the way: coded syntax and substitution, symbolic unfolding, internal
    strong induction ---\\
    \key{evidence that RGA supports nontrivial mathematical reasoning.}
  \end{center}
\end{frame}

\begin{frame}{5 $\cdot$ \emph{Idealizing Useful Fictions in OGA}}
  {\small Ford, arXiv:2608.17862 (Aug 2026)}
  \vfill
  RGA can quantify over its own computations --- but \key{cannot certify that its own
  unbounded searches have definite yes-or-no answers}.
  \vfill
  OGA closes that openness with \key{exactly one rule}:
  \begin{center}
    \key{ATI}, the $\omega$-grounded universal\\[0.3em]
    \small\emph{if every numeric instance of a universal is certified decided,\\
    the universal is certified decided.}
  \end{center}
  \vfill
  \small
  Otherwise identical to RGA, symbol for symbol. Every consequence machine-checked.
  \vfill
  \begin{center}
    Certificates become \key{abundant}: every totality question about a computable
    function is certified to \emph{have} an answer ---\\
    whether or not anyone can produce it.
  \end{center}
\end{frame}

\begin{frame}{What that one rule costs, and buys}
  \small
  \begin{center}
  \begin{tabular}{@{}lll@{}}
    \toprule
                        & \textbf{RGA} & \textbf{OGA} \\
    \midrule
    a certificate says  & a computation settles it
                        & the question \emph{has} an answer \\
    provability         & r.e. & r.e., primitive-recursive checker \\
    truth               & \key{r.e.} & \key{deliberately \emph{not} r.e.} \\
    the residue         & $\omega$-incompleteness & \key{fictions} \\
    \bottomrule
  \end{tabular}
  \end{center}
  \vfill
  \normalsize
  In OGA the G\"odel sentence is classified \emph{unconditionally} as a
  \key{genuine fiction}: neither provable nor refutable, yet \emph{valued} ---
  and carrying a \key{computable pedigree} recording exactly what adopting it
  as an axiom commits you to.
  \vfill
  \begin{center}\small
    Adoption is itself a theorem suite: extending OGA by any finite stock of true
    fictions is consistent,\\ and independently certified adoptions can never collide.
  \end{center}
\end{frame}

\begin{frame}{6 $\cdot$ \emph{Fact and Fiction in Grounded Set Theory}}
  {\small Ford (Aug 2026)}
  \vfill
  \begin{center}\small
    ``Classical set theory buys its power on credit:\\
    its axioms assert totalities no computation could survey.''
  \end{center}
  \vfill
  Sets are \key{formula codes}; membership is \emph{defined} via provability ---
  no primitive $\in$, and no set axioms assumed.
  \vfill
  \small
  \begin{center}
  \begin{tabular}{@{}ll@{}}
    \toprule
    theorems & empty, pairing, separation, foundation, choice, complement,
               \key{Tarski universe} \\
    fail     & \bad{union, replacement, powerset} --- the \emph{collecting} axioms,
               recursion-theoretically \\
    graded   & \amber{extensionality} --- pointwise provable, uniformly a G\"odel sentence \\
    \bottomrule
  \end{tabular}
  \end{center}
  \vfill
  \begin{center}
    ``ZF fails computably, but holds fictionally,\\
    \key{and the fiction is measurable}.''
  \end{center}
\end{frame}

\begin{frame}{What the programme has, and has not}
  \vfill
  \begin{columns}[T]
    \begin{column}{0.47\textwidth}
      \underline{Settled}\\[0.8em]
      \small
      Unrestricted recursive definitions, admitted with \emph{no} obligation at
      definition time\\[0.5em]
      Consistency --- machine-checked\\[0.5em]
      Semantic completeness\\[0.5em]
      Quantifiers\\[0.5em]
      Its own truth predicate\\[0.5em]
      Incompleteness \emph{priced} --- fictions with pedigrees\\[0.5em]
      A graded set theory
    \end{column}
    \begin{column}{0.47\textwidth}
      \underline{Open}\\[0.8em]
      \small
      Proof-theoretic strength\\[0.5em]
      Whether \hq\ obligations can be discharged \emph{cheaply enough}\\[0.5em]
      Datatypes beyond $\mathbb{N}$\\[0.5em]
      Notation, tactics, automation\\[0.5em]
      \key{Whether anyone can reason in it}
    \end{column}
  \end{columns}
  \vfill
  \begin{center}
    \small Church and Curry were defeated by the \emph{consistency} question.\\
    \key{That one is answered. The usability question is not.}
  \end{center}
\end{frame}

% =====================================================================
\section{This work}
% =====================================================================

\begin{frame}{The gap}
  \vfill
  \begin{center}\Large
    A logic is not a tool.
  \end{center}
  \vfill
  \begin{itemize}\itemsep0.8em
    \item The metatheory is done in \key{Isabelle/HOL} --- which adds HOL's axioms
          to GD's trusted base, and embeds GD's primitives inside a very different logic
    \item Good for \emph{studying} GA. Wrong for \emph{reasoning in} GA.
    \item A foundational system should sit atop a \key{minimal} framework
  \end{itemize}
  \vfill
  \begin{center}\small
    Kehrli (2025): formalize GA atop \key{Isabelle/Pure} instead.
  \end{center}
  \vfill
\end{frame}

\begin{frame}{What makes GA hard to use}
  \vfill
  In classical logic, $p\ \mathsf{B}$ is free --- it is the law of excluded middle.
  \vfill
  \begin{center}\Large
    In GA it is \key{work}.
  \end{center}
  \vfill
  Every \hq\ premise is a proof obligation that a human or a tactic must discharge.
  \vfill
  \pause
  \begin{center}
    \large So the engineering question is sharp and checkable:\\[0.6em]
    \key{can these obligations be discharged automatically often enough\\
    that ordinary reasoning stays tolerable?}
  \end{center}
  \vfill
\end{frame}

\begin{frame}{This work}
  \vfill
  Building a usable theorem prover for grounded deduction on Isabelle/Pure.
  \vfill
  \begin{itemize}\itemsep0.8em
    \item Proof search and tactics for \hq\ obligations
    \item Conditional rewriting where equality is not reflexive
    \item Induction and case analysis under grounding premises
    \item Encoding inductive datatypes in a system with only $\mathbb{N}$
  \end{itemize}
  \vfill
  \begin{center}\small
    \emph{[ details follow ]}
  \end{center}
  \vfill
\end{frame}

\begin{frame}[plain]
  \vfill
  \begin{center}
    {\Large Reasoning Around Recursion}\\[2em]
    \small
    Curry proved you cannot have both general recursion\\
    and ordinary implication.\\[1em]
    Everyone since has given up recursion.\\[1em]
    \key{Grounded arithmetic gives up something in implication instead.}
  \end{center}
  \vfill
\end{frame}

\end{document}
